% Capítulo 2: Objetivos del TFM

\chapter{Objetivos del TFM}

\label{Chapter2}

\lhead{Capítulo 2. \emph{Objetivos}}

\textit{[En un párrafo explica brevemente los contenidos de este capítulo. Se recomienda un tamaño de 1 a 3 páginas para este capítulo.]}

%----------------------------------------------------------------------------------------

\section{Objetivo general}

\textit{[El objetivo general es lo que pretendes hacer con el TFM: el último paso a realizar. Redáctalo en infinitivo y con un solo verbo principal.]}

%----------------------------------------------------------------------------------------

\section{Objetivos específicos}

\textit{[Los objetivos específicos son las partes en las que subdivides el trabajo para completar el objetivo general.]}

\begin{itemize}
    \item \textit{[No uses más de 4--6 objetivos.]}
    \item \textit{[Siempre redactados en infinitivo.]}
    \item \textit{[Los objetivos no son deseos para el futuro: cuando termines el proyecto se deben haber cumplido todos.]}
    \item \textit{[Lo mejor es poner un objetivo para el marco teórico, otro para el metodológico, otro para el empírico, etc. Es decir, trocear el objetivo final en pequeñas metas volantes.]}
\end{itemize}
